\documentclass[a4paper,11pt]{article}
\usepackage{lmodern}
\usepackage[T1]{fontenc}
\usepackage[utf8]{inputenc}
\usepackage{amsmath,amssymb,amsthm}
\usepackage{mathtools}
\usepackage{booktabs}
\usepackage{longtable}
\usepackage{array}
\usepackage{hyperref}
\usepackage{graphicx}
\usepackage{float}
\usepackage{geometry}
\geometry{margin=25mm}

\hypersetup{
  colorlinks=true,
  linkcolor=blue,
  urlcolor=blue,
  citecolor=blue
}

\theoremstyle{definition}
\newtheorem{theorem}{Theorem}[section]
\newtheorem{proposition}[theorem]{Proposition}
\newtheorem{lemma}[theorem]{Lemma}
\newtheorem{corollary}[theorem]{Corollary}
\newtheorem{definition}[theorem]{Definition}
\newtheorem{remark}[theorem]{Remark}
\newtheorem{observation}[theorem]{Observation}

\setlength{\parindent}{1em}
\setlength{\parskip}{0.5em}

\providecommand{\tightlist}{\setlength{\itemsep}{0pt}\setlength{\parskip}{0pt}}
\providecommand{\real}[1]{#1}
\begin{document}

\section{What Is the Diameter of the Circumscribed Hypersphere of a
Four-Dimensional Hyperrectangle with Unit
Volume?}\label{what-is-the-diameter-of-the-circumscribed-hypersphere-of-a-four-dimensional-hyperrectangle-with-unit-volume}

\subsection{── Spacetime Dimensionality, the Origin of Time, and the
Beginning of the Universe as Derived from a Discrete Packing Model
──}\label{spacetime-dimensionality-the-origin-of-time-and-the-beginning-of-the-universe-as-derived-from-a-discrete-packing-model}

\begin{center}\rule{0.5\linewidth}{0.5pt}\end{center}

\textbf{Author}: Noriaki Kihara (WF System Co., Ltd.)

\textbf{Date}: April 2026 (Rev.2)

\textbf{DOI:}
\href{https://doi.org/10.5281/zenodo.19834940}{10.5281/zenodo.19834940}
(v2; v1:
\href{https://doi.org/10.5281/zenodo.19533313}{10.5281/zenodo.19533313})

\begin{center}\rule{0.5\linewidth}{0.5pt}\end{center}

\subsection{0. Scope of This Paper's Claims (Required
Reading)}\label{scope-of-this-papers-claims-required-reading}

This paper is a \textbf{purely mathematical thought experiment} and
makes no claims about the actual universe. To prevent misreading, the
paper's \textbf{claims} and \textbf{non-claims} are stated explicitly at
the outset.

\subsubsection{What This Paper Claims}\label{what-this-paper-claims}

\begin{enumerate}
\def\labelenumi{\arabic{enumi}.}
\tightlist
\item
  Starting from two assumptions --- ``discreteness'' and ``stability''
  --- together with a ``volume conservation law,'' certain geometric and
  combinatorial consequences are \textbf{deductively derived}.
\item
  Among those consequences is the mathematical fact that
  \textbf{\(n = 4\)} is uniquely selected as the smallest non-trivial
  dimension for which \(\sqrt{n}\) is an integer.
\item
  For the sign-reversal operation on signed volumes in four dimensions,
  there exists a minimally economical option --- \textbf{single-axis
  reversal} --- which produces the structural asymmetry of ``1 versus
  3.''
\item
  The mathematical structures above possess forms \textbf{structurally
  analogous} to physical concepts such as four-dimensional spacetime,
  the Minkowski metric, and the beginning of the universe.
\item
  The argument in this paper is essentially \textbf{combinatorics of
  information}: the ``zero-dimensional state'' at the starting point
  refers not to a physical point but to \textbf{structural information}
  held statically without spatial degrees of freedom (see the
  supplementary note in Section 8 for details).
\end{enumerate}

\subsubsection{What This Paper Does Not
Claim}\label{what-this-paper-does-not-claim}

\begin{enumerate}
\def\labelenumi{\arabic{enumi}.}
\tightlist
\item
  \textbf{This paper does not explain why the actual universe is
  four-dimensional.} Why physical spacetime has four dimensions is an
  unsolved problem in physics; this paper makes no claim whatsoever that
  its model provides the answer.
\item
  \textbf{This paper is not a physical theory.} No rigorous
  correspondence has been established with general relativity, quantum
  field theory, the Standard Model, string theory, quantum gravity, or
  related frameworks. Nor does this paper claim that its
  ``hyperrectangles'' correspond to physical entities.
\item
  \textbf{Terms such as ``time,'' ``space,'' ``matter,'' ``antimatter,''
  and ``Big Bang'' as used in this paper are metaphors.} They should not
  be identified with the same-named concepts in physics. They are
  borrowed terms used solely to describe features of mathematical
  structures.
\item
  \textbf{This paper makes no empirical predictions.} It provides no
  quantitative predictions verifiable by observation or experiment, and
  constitutes neither an alternative to, nor a refutation or
  corroboration of, any existing physical theory.
\item
  \textbf{This paper makes no cosmological interpretations whatsoever.}
  Cosmological language appearing in this paper (e.g., ``the beginning
  of the universe,'' ``emergence of time,'' ``universe and
  anti-universe'') is entirely \textbf{metaphorical language} intended
  to convey features of mathematical structures intuitively; it is not a
  hypothesis, claim, or prediction regarding the actual universe.
\end{enumerate}

\subsubsection{The Position of This
Paper}\label{the-position-of-this-paper}

This paper is a record of a thought experiment conducted interactively,
beginning from the elementary geometric question of the circumscribed
sphere diameter of a four-dimensional hyperrectangle and exploring how
far interesting structures can be derived from a very small number of
assumptions. It is \textbf{mathematical play} and \textbf{intellectual
exercise}. Readers are invited to read it not as a physical claim but as
an \textbf{illustration} of how concepts such as symmetry, integrality,
and conservation laws interact with one another.

Discussing the correspondence between mathematical structures and
physical reality requires rigorous formalization, connection to existing
physics, and observational verification far beyond the scope of this
paper. Those tasks lie outside its reach.

\subsubsection{About the Revised Version
(Rev.2)}\label{about-the-revised-version-rev.2}

The original §4 treated ``a sphere circumscribing a hypercube of side
length \((2k+1)\).'' However, this merely placed a cube-shaped region
inside a sphere and did not pack the sphere itself (the packing density
was permanently capped at \(2/\pi^2 \approx 0.2026\)). The revised
version reformulates the original question as ``densely packing unit
cubes into a sphere of radius \(R = 2k+1\),'' and derives a sequence
\(N(k)\) whose packing density asymptotically converges to 1 as
\(k \to \infty\). The arguments from §7 onwards regarding signed volume,
the time axis, and cosmological interpretation are preserved as they
stand, with the sequence \((2k+1)^4\) replaced by the new packing count
\(N(k)\).

\begin{center}\rule{0.5\linewidth}{0.5pt}\end{center}

\subsection{1. A Naive Question: The Diameter of the Circumscribed
Hypersphere}\label{a-naive-question-the-diameter-of-the-circumscribed-hypersphere}

Consider a four-dimensional hypercube with side length 1. What is the
diameter of the four-dimensional hypersphere circumscribed about this
hypercube --- that is, the hypersphere passing through all of its
vertices?

The length of the space diagonal of a four-dimensional hypercube is

\[d = \sqrt{1^2 + 1^2 + 1^2 + 1^2} = \sqrt{4} = 2\]

Since the diameter of the circumscribed hypersphere equals this
diagonal, the \textbf{diameter is 2} (radius 1).

\begin{center}\rule{0.5\linewidth}{0.5pt}\end{center}

\subsection{2. Symmetric Packing: How Many Can Be
Packed?}\label{symmetric-packing-how-many-can-be-packed}

Next, consider packing multiple copies of this four-dimensional
hypercube inside the circumscribed hypersphere while preserving
rotational symmetry.

Arranging unit hypercubes symmetrically around the origin, each
hypercube can occupy one ``orthant'' of four-dimensional space. Since
there are \(2^4 = 16\) orthants in four dimensions, \textbf{16} unit
hypercubes share a corner at the origin and pack together without gaps,
forming a hypercube \([-1,1]^4\) with side length 2.

This arrangement preserves the full symmetry group \(B_4\) of the
hypercube (order 384) and is the most symmetric packing possible. The
radius of the circumscribed hypersphere --- the distance from the center
to a vertex --- is \(\sqrt{1^2+1^2+1^2+1^2} = 2\), giving a
\textbf{diameter of 4}.

A comparison across dimensions is shown below:

{\def\LTcaptype{none} % do not increment counter
\begin{longtable}[]{@{}
  >{\centering\arraybackslash}p{(\linewidth - 4\tabcolsep) * \real{0.3333}}
  >{\centering\arraybackslash}p{(\linewidth - 4\tabcolsep) * \real{0.3333}}
  >{\centering\arraybackslash}p{(\linewidth - 4\tabcolsep) * \real{0.3333}}@{}}
\toprule\noalign{}
\begin{minipage}[b]{\linewidth}\centering
Dimension \(n\)
\end{minipage} & \begin{minipage}[b]{\linewidth}\centering
Count \(2^n\)
\end{minipage} & \begin{minipage}[b]{\linewidth}\centering
Circumscribed sphere diameter \(2\sqrt{n}\)
\end{minipage} \\
\midrule\noalign{}
\endhead
\bottomrule\noalign{}
\endlastfoot
2 & 4 & \(2\sqrt{2} \approx 2.83\) \\
3 & 8 & \(2\sqrt{3} \approx 3.46\) \\
\textbf{4} & \textbf{16} & \textbf{4 (integer)} \\
\end{longtable}
}

\begin{center}\rule{0.5\linewidth}{0.5pt}\end{center}

\subsection{3. The Special Nature of Four
Dimensions}\label{the-special-nature-of-four-dimensions}

A noteworthy fact emerges here. The circumscribed sphere diameter
\(2\sqrt{n}\) is an integer only when \(\sqrt{n}\) is an integer ---
that is, only when \(n\) is a perfect square.

\[n = 1, 4, 9, 16, 25, \ldots\]

Excluding the trivial case \(n = 1\) (a one-dimensional line segment),
\textbf{\(n = 4\) is the smallest non-trivial dimension satisfying
integrality}. This is not a coincidence; it is an essential property
that underpins the discussion that follows.

\begin{center}\rule{0.5\linewidth}{0.5pt}\end{center}

\subsection{4. Formalization of Dense
Packing}\label{formalization-of-dense-packing}

\subsubsection{4.1 Setup}\label{setup}

Consider the four-dimensional ball centered at the origin

\[B(R) := \{x \in \mathbb{R}^4 : \|x\|_2 \le R\}, \quad R = 2k+1 \;\; (k \in \mathbb{Z}_{\ge 0})\]

Pack unit hypercubes whose centers lie at integer lattice points
\((x, y, z, w) \in \mathbb{Z}^4\) inside \(B(R)\), requiring that they
not protrude. Each unit cube must be \textbf{fully contained} in
\(B(R)\).

\subsubsection{4.2 Packing Condition}\label{packing-condition}

Of the 16 vertices of the unit cube centered at \((x, y, z, w)\), the
one farthest from the origin is

\[P_{\max} = \left(|x|+\tfrac{1}{2},\ |y|+\tfrac{1}{2},\ |z|+\tfrac{1}{2},\ |w|+\tfrac{1}{2}\right)\]

The necessary and sufficient condition for the cube to be fully
contained in \(B(R)\) is \(\|P_{\max}\|_2 \le R\), i.e.,

\[\boxed{\;\left(|x|+\tfrac{1}{2}\right)^2 + \left(|y|+\tfrac{1}{2}\right)^2 + \left(|z|+\tfrac{1}{2}\right)^2 + \left(|w|+\tfrac{1}{2}\right)^2 \;\le\; (2k+1)^2\;}\]

We \textbf{define} \(N(k)\) as the number of integer tuples
\((x,y,z,w) \in \mathbb{Z}^4\) satisfying this inequality.

\subsubsection{\texorpdfstring{4.3 The Sequence
\(N(k)\)}{4.3 The Sequence N(k)}}\label{the-sequence-nk}

By direct computation:

{\def\LTcaptype{none} % do not increment counter
\begin{longtable}[]{@{}cccc@{}}
\toprule\noalign{}
\(k\) & \(R = 2k+1\) & \(N(k)\) & Shell increment
\(a_k = N(k) - N(k-1)\) \\
\midrule\noalign{}
\endhead
\bottomrule\noalign{}
\endlastfoot
0 & 1 & 1 & --- \\
1 & 3 & 137 & 136 \\
2 & 5 & 1,545 & 1,408 \\
3 & 7 & 7,281 & 5,736 \\
4 & 9 & 22,409 & 15,128 \\
5 & 11 & 53,161 & 30,752 \\
6 & 13 & 108,081 & 54,920 \\
7 & 15 & 199,953 & 91,872 \\
8 & 17 & 337,417 & 137,464 \\
9 & 19 & 537,409 & 199,992 \\
10 & 21 & 818,145 & 280,736 \\
\end{longtable}
}

The sequence \(N(k)\) is not the fourth-power sequence \((2k+1)^4\) of
the original version, but a number-theoretic sequence based on lattice
point counting inside a four-dimensional ball (see §5).

\subsubsection{4.4 Packing Density}\label{packing-density}

The volume of a four-dimensional ball is \(V_4(R) = (\pi^2/2) R^4\). The
packing density \(\rho(k) := N(k) / V_4(2k+1)\) is:

{\def\LTcaptype{none} % do not increment counter
\begin{longtable}[]{@{}ccc@{}}
\toprule\noalign{}
\(k\) & \(V_4(R)\) & \(\rho(k)\) \\
\midrule\noalign{}
\endhead
\bottomrule\noalign{}
\endlastfoot
0 & 4.93 & 0.2026 \\
1 & 399.72 & 0.3427 \\
2 & 3,084.25 & 0.5009 \\
5 & 72,250.44 & 0.7358 \\
10 & 959,725.27 & 0.8525 \\
\end{longtable}
}

\textbf{Proposition 4.1}:
\(\displaystyle \lim_{k \to \infty} \rho(k) = 1\).

\textbf{Sketch of proof}: The number of integer lattice points
satisfying the packing condition is asymptotic to the volume of the
region
\(\Omega_R := \{x \in \mathbb{R}^4 : \sum_{i=1}^{4} (|x_i|+1/2)^2 \le R^2\}\).
\(\Omega_R\) is \(B(R)\) shrunk inward by half a unit width, and its
volume tends to \(V_4(R)\) as \(R \to \infty\). The error term is the
surface-area term \(O(R^3)\), giving \(\rho(k) = 1 + O(R^{-1}) \to 1\).
\(\square\)

In the original version the packing density was the constant
\(2/\pi^2 \approx 0.2026\) independent of \(k\), but in the revised
version it grows asymptotically as \(0.2026 \to 1\) as \(k\) increases.
This is consistent with the intuition that ``as the ball grows, the gaps
are filled in by unit cubes.''

\begin{center}\rule{0.5\linewidth}{0.5pt}\end{center}

\subsection{5. The Special Nature of Four Dimensions (Number-Theoretic
Background)}\label{the-special-nature-of-four-dimensions-number-theoretic-background}

\subsubsection{5.1 Integrality of the Diagonal
Solution}\label{integrality-of-the-diagonal-solution}

Taking the symmetric (diagonal) solution \(x = y = z = w\) in four
dimensions,

\[4x^2 = R^2 \quad\Longleftrightarrow\quad R = 2|x|\]

holds. This is a direct consequence of the fact that \(\sqrt{4} = 2\) is
an integer, and does not hold in other low dimensions.

In the context of the packing problem, taking the half-integer
\(|x| + 1/2 = (2k+1)/2\) for \(R = 2k+1\) (odd) gives

\[4 \cdot \left(\frac{2k+1}{2}\right)^2 = (2k+1)^2 = R^2\]

\subsubsection{5.2 Spherical Contact of the Corner
Cubes}\label{spherical-contact-of-the-corner-cubes}

\textbf{Proposition 5.1}: For every \(k \ge 0\), the unit cubes centered
at \((\pm k, \pm k, \pm k, \pm k)\) --- one in each orthant,
\(2^4 = 16\) ``corner cubes'' in total --- have their farthest vertices
lying exactly on the sphere \(\partial B(2k+1)\).

\textbf{Proof}: The farthest vertex of the cube centered at
\((\pm k)^4\) is \((\pm(k+1/2))^4\). Its distance from the origin is

\[\sqrt{4 \cdot (k+1/2)^2} \;=\; 2(k+1/2) \;=\; 2k+1 \;=\; R \qquad \square\]

This is a phenomenon specific to four dimensions: the vertices
(half-integer-coordinate lattice points) of unit cubes lie at integer
distances on the circumscribed sphere. This does not hold in three
dimensions because \(\sqrt{3}\) is irrational.

\subsubsection{5.3 Lagrange--Jacobi Number-Theoretic
Background}\label{lagrangejacobi-number-theoretic-background}

Performing the change of variables \(p_i := 2|x_i| + 1\) (positive odd
integers) in the packing condition and multiplying both sides by 4 gives

\[p_1^2 + p_2^2 + p_3^2 + p_4^2 \;\le\; (4k+2)^2\]

The boundary equality \(p_1^2 + \cdots + p_4^2 = (4k+2)^2\) reduces to
the number-theoretic problem of representations as sums of four squares.

\textbf{Lagrange's four-square theorem (1770)}: Every non-negative
integer can be expressed as the sum of at most four squares. This means
that four dimensions is the smallest dimension in which ``every integer
is realized by lattice points in \(n\) dimensions.''

\textbf{Jacobi's four-square theorem (1834)}: The number \(r_4(N)\) of
ways to express \(N\) as the sum of four integer squares has the closed
form

\[r_4(N) = \begin{cases} 8\,\sigma(N) & (4 \nmid N) \\ 24\,\sigma(N_{\text{odd}}) & (4 \mid N) \end{cases}\]

(where \(\sigma\) is the divisor sum and \(N_{\text{odd}}\) is the odd
part of \(N\)). This simple closed form is specific to four dimensions;
in three dimensions, \(r_3(N)\) is far more complex, involving class
numbers.

The packing count \(N(k)\) is directly tied to this number-theoretic
structure.

\begin{center}\rule{0.5\linewidth}{0.5pt}\end{center}

\subsection{6. Four Dimensions as a Stability Condition
(Reformulated)}\label{four-dimensions-as-a-stability-condition-reformulated}

The stability condition of the original §6 --- ``symmetric packing
within the circumscribed hypersphere is possible and the diameter is an
integer'' --- is reformulated in this revised version as follows.

\textbf{Stability Condition (Revised)}: A dimension \(n\) is called
\textbf{stable} if it satisfies all of the following:

\begin{enumerate}
\def\labelenumi{\arabic{enumi}.}
\tightlist
\item
  Dense packing of unit cubes into a ball of radius
  \(R \in \mathbb{Z}_{>0}\) is well-defined.
\item
  The packing density asymptotically converges to 1 as \(R \to \infty\).
\item
  The vertices of the ``corner cubes'' lie on integer-distance lattice
  points on the sphere (integrality of \(\sqrt{n}\)).
\item
  The representation count \(r_n(N)\) for sums of \(n\) squares has a
  simple closed form.
\end{enumerate}

The smallest non-trivial dimension satisfying all of these is
\textbf{\(n = 4\)}.

\begin{itemize}
\tightlist
\item
  For \(n = 3\), conditions (3) and (4) fail. Since \(\sqrt{3}\) is
  irrational, the corner cubes do not lie at integer-distance lattice
  points on the sphere, and \(r_3(N)\) also follows a complex formula
  involving class numbers. The symmetry group \(B_3\) (order 48) is also
  weaker than \(B_4\) (order 384).
\item
  For \(n = 9, 16, \ldots\), condition (3) does hold, but \(n = 4\) is
  selected by minimality.
\end{itemize}

\begin{center}\rule{0.5\linewidth}{0.5pt}\end{center}

\subsection{7. Signed Volume: Negative
Volume}\label{signed-volume-negative-volume}

At this point, the direction of the discussion turns, and a
\textbf{sign} is introduced for the volume of a hyperrectangle.

Ordinary geometric volume is always non-negative, but in linear algebra
one can consider signed volume --- oriented volume defined by a
determinant. Regarding a four-dimensional hyperrectangle as a
parallelepiped spanned by four vectors
\(\mathbf{v}_1, \mathbf{v}_2, \mathbf{v}_3, \mathbf{v}_4\), its signed
volume is
\(\det[\mathbf{v}_1 \; \mathbf{v}_2 \; \mathbf{v}_3 \; \mathbf{v}_4]\).

Reversing \(k\) vectors (multiplying by \(-1\)) multiplies the sign by
\((-1)^k\). Therefore:

\begin{itemize}
\tightlist
\item
  Reversing 1 vector: \((-1)^1 = -1\) (volume becomes negative)
\item
  Reversing 2 vectors: \((-1)^2 = +1\) (remains positive)
\item
  Reversing 3 vectors: \((-1)^3 = -1\) (becomes negative)
\item
  Reversing 4 vectors: \((-1)^4 = +1\) (returns to positive)
\end{itemize}

Negative volume is realized by reversing exactly 1 or exactly 3 axes.

\begin{center}\rule{0.5\linewidth}{0.5pt}\end{center}

\subsection{8. Initial Conditions: A Zero-Volume Compressed
State}\label{initial-conditions-a-zero-volume-compressed-state}

From here, construction of the cosmological model begins. The following
initial conditions are assumed.

It is assumed that \textbf{the beginning of the universe started from a
zero-dimensional state} --- a maximally compressed state with no
dimensional degrees of freedom. The following are packed into this
state:

\begin{itemize}
\tightlist
\item
  \(N(k)\) four-dimensional hyperrectangles with volume \(+1\)
\item
  \(N(k)\) four-dimensional hyperrectangles with volume \(-1\)
\item
  \textbf{The total signed volume is zero}
\end{itemize}

The whole is compressed inside a circumscribed hypersphere of radius 1.
This is an extreme degenerate state in which \(2N(k)\) hyperrectangles
--- the information of a ball region of volume \(V_4(2k+1)\) that would
be packed at density \(\rho(k)\) --- are superposed at a single point of
zero degrees of freedom.

Although positive and negative contributions cancel to give total volume
zero, the information of ``count'' and ``symmetric structure'' is
retained in the number \(N(k)\).

\subsubsection{Supplementary Note: The Meaning of ``Information'' in
Zero
Dimensions}\label{supplementary-note-the-meaning-of-information-in-zero-dimensions}

The natural question arises: ``How can \(N(k)\) hyperrectangles exist in
zero dimensions?'' This paper's position on this question is clarified
below.

\textbf{Zero-dimensional} in this paper refers to a state with
absolutely no spatial extent (degrees of freedom). However, being
zero-dimensional does not mean having no information. Zero dimensions
can carry information \(I\) of the kind ``present or absent.''

The ``information'' referred to here is not limited to representations
such as bit strings of 1s and 0s. It may be real-valued, complex-valued,
matrix-valued, tensor-valued, or any combination thereof. The only
requirement is simply:

\begin{quote}
\textbf{What information exists or does not exist is expressed
statically without any spatial extent.}
\end{quote}

The description at the starting point of this paper --- ``\(N(k)\)
positive hyperrectangles and \(N(k)\) negative hyperrectangles'' ---
does not refer to objects arranged in space, but rather describes
\textbf{the structure of this static information \(I\)}. It should be
read as: combinatorial information concerning count, sign, and symmetry
is held as a state with zero degrees of freedom.

In other words, the zero-dimensional state in this paper is not a
``physical situation in which objects are packed into a point,'' but
rather \textbf{``structural information prior to its unfolding.''} The
release of degrees of freedom corresponds to the process in which this
information is \textbf{read out} in a dimensioned form.

In this sense, the argument in this paper is \textbf{combinatorics of
information} before it is a physical picture, and the assertion that
``information exists in zero dimensions'' is mathematically
non-contradictory.

\begin{center}\rule{0.5\linewidth}{0.5pt}\end{center}

At this point, the most important principle is stated explicitly:

\begin{quote}
\textbf{The volume conservation law cannot be violated.}
\end{quote}

If negative volume were not permitted, the beginning of the universe
would require positive volume to emerge from zero volume, violating the
conservation law. \textbf{Negative volume is not a strange assumption;
it is a logical necessity for preserving the conservation law.}

\begin{center}\rule{0.5\linewidth}{0.5pt}\end{center}

\subsection{9. The First Degree of Freedom: Unfolding into One
Dimension}\label{the-first-degree-of-freedom-unfolding-into-one-dimension}

Next, suppose that one-dimensional degrees of freedom are suddenly
granted. At this point, the following two rules are introduced:

\begin{enumerate}
\def\labelenumi{\arabic{enumi}.}
\tightlist
\item
  \textbf{Minimize the packing density} (spread as thinly as possible).
\item
  \textbf{Minimize the total four-dimensional hypervolume} (minimize the
  spread in four-dimensional directions).
\end{enumerate}

Satisfying both rules simultaneously, the limiting solution is for
positive--negative pairs to be arranged point-by-point along a
one-dimensional axis. The four-dimensional cross-section of each pair
approaches zero, and the structure extends only in the length direction
by \(N(k)\).

As a result, the whole degenerates into ``a one-dimensional line of zero
thickness and length \(N(k)\).'' This is a transition that abandons
four-dimensional degrees of freedom in favor of one-dimensional degrees
of freedom; the information of the original four-dimensional structure
remains fully encoded in the one-dimensional quantity \(N(k)\) --- the
length.

\begin{center}\rule{0.5\linewidth}{0.5pt}\end{center}

\subsection{10. Unfolding into Two Dimensions and Diameter
Minimization}\label{unfolding-into-two-dimensions-and-diameter-minimization}

When the degrees of freedom become two-dimensional, a third rule is
added:

\begin{enumerate}
\def\labelenumi{\arabic{enumi}.}
\setcounter{enumi}{2}
\tightlist
\item
  \textbf{Minimize the diameter of the circumscribed hypersphere.}
\end{enumerate}

In one dimension, the length \(N(k)\) directly becomes the diameter,
which is excessively large. With two-dimensional degrees of freedom, the
diameter can be minimized by \textbf{folding this into a square
configuration}.

Arranging \(N(k)\) elements in a square, each side has length

\[\sqrt{N(k)} = N(k)^{1/2}\]

and the diameter of the circumscribed hypersphere shrinks to the
square's diagonal:

\[\sqrt{2} \cdot N(k)^{1/2}\]

The four-dimensional directions remain collapsed to a point; the
structure is unfolded in the two-dimensional direction as an
\(N(k)^{1/2} \times N(k)^{1/2}\) lattice.

\begin{center}\rule{0.5\linewidth}{0.5pt}\end{center}

\subsection{11. Pattern of Side Lengths Across
Dimensions}\label{pattern-of-side-lengths-across-dimensions}

A pattern becomes visible. Unfolding \(N(k)\) elements into \(n\)
degrees of freedom gives a side length of \(N(k)^{1/n}\). Using the
asymptotics \(N(k) \sim (\pi^2/2) R^4\) (with \(R = 2k+1\),
\(k \to \infty\)):

{\def\LTcaptype{none} % do not increment counter
\begin{longtable}[]{@{}
  >{\centering\arraybackslash}p{(\linewidth - 4\tabcolsep) * \real{0.3333}}
  >{\centering\arraybackslash}p{(\linewidth - 4\tabcolsep) * \real{0.3333}}
  >{\centering\arraybackslash}p{(\linewidth - 4\tabcolsep) * \real{0.3333}}@{}}
\toprule\noalign{}
\begin{minipage}[b]{\linewidth}\centering
Degrees of freedom \(n\)
\end{minipage} & \begin{minipage}[b]{\linewidth}\centering
Side length
\end{minipage} & \begin{minipage}[b]{\linewidth}\centering
Asymptotic form (\(k \to \infty\))
\end{minipage} \\
\midrule\noalign{}
\endhead
\bottomrule\noalign{}
\endlastfoot
1 & \(N(k)\) & \((\pi^2/2)\,(2k+1)^4\) \\
2 & \(N(k)^{1/2}\) & \((\pi^2/2)^{1/2}\,(2k+1)^2\) \\
3 & \(N(k)^{1/3}\) & \((\pi^2/2)^{1/3}\,(2k+1)^{4/3}\) \\
\textbf{4} & \(N(k)^{1/4}\) & \((\pi^2/2)^{1/4}\,(2k+1)\) \\
\end{longtable}
}

At \(n = 4\), the asymptotic form scales \textbf{linearly} with the ball
radius \(R = 2k+1\) (with coefficient
\((\pi^2/2)^{1/4} \approx 1.4894\)). This means the unit-cube lattice
scale and the ball radius asymptotically coincide, and is reinterpreted
as the \textbf{asymptotic recovery} of the symmetric four-dimensional
hypercube structure. At \(n = 3\), the asymptotic form is
\((2k+1)^{4/3}\), which diverges faster than the radius and breaks
dimensional consistency. For all \(n < 4\), the side length and the
radius do not match in scale.

\begin{center}\rule{0.5\linewidth}{0.5pt}\end{center}

\subsection{12. Why the System Settles in Four
Dimensions}\label{why-the-system-settles-in-four-dimensions}

Searching for a ``dimension in which the system can settle'' --- one
satisfying all the rules (maintenance of symmetry, spherical contact of
corner cubes via integrality of \(\sqrt{n}\), minimization of
circumscribed sphere diameter, minimization of hypervolume)
simultaneously --- one finds:

\begin{itemize}
\tightlist
\item
  \textbf{\(n = 0\)}: The initial state. Total volume zero, compressed
  into a sphere of radius 1. The starting point.
\item
  \textbf{\(n = 1\)}: A line of length \(N(k)\). The diameter explodes
  in size, violating diameter minimization. A transient state.
\item
  \textbf{\(n = 2\)}: \(\sqrt{n}\) is irrational, so corner cubes cannot
  exactly inscribe on the sphere. An intermediate state.
\item
  \textbf{\(n = 3\)}: \(\sqrt{n}\) is irrational. The side length
  \(N(k)^{1/3}\) does not match the scale of the ball radius \(R\).
  Unstable.
\item
  \textbf{\(n = 4\)}: \(\sqrt{4} = 2\), so corner cubes inscribe exactly
  on the sphere (Proposition 5.1). Full recovery of \(B_4\) symmetry.
  The side length asymptotes to \((\pi^2/2)^{1/4}(2k+1)\), matching the
  scale of the ball radius \(R = 2k+1\). The system is stable with
  diameter \(2(2k+1)\). \textbf{The first true equilibrium point.}
\end{itemize}

Further equilibrium points exist at \(n = 9, 16, \ldots\), but the
system stops at the minimum solution of 4, balancing diameter
minimization with the preference not to increase dimensionality.

\begin{quote}
\textbf{Starting from zero dimensions and releasing degrees of freedom
one by one, none of the dimensions 1 through 3 can sustain a stable
equilibrium; only upon reaching four dimensions are all rules
simultaneously satisfied and the system settles.}
\end{quote}

Viewing this sequence, after 0 the next dimension at which the system
can settle is 4; given degrees of freedom, the system ends up stable in
four dimensions. Four-dimensionality is not an assumption but a
consequence.

\begin{center}\rule{0.5\linewidth}{0.5pt}\end{center}

\subsection{13. The Origin of the Time Axis: Why Exactly One Axis Is
Special}\label{the-origin-of-the-time-axis-why-exactly-one-axis-is-special}

Realizing negative volume requires reversing an odd number of axes. The
most economical operation is \textbf{reversing exactly one axis},
producing \((-1)^1 = -1\).

By reversing \textbf{exactly one} of the four axes --- distinguishing it
from the others --- negative volume is realized. The remaining three
axes are not involved in the reversal and remain on equal footing.

\begin{quote}
\textbf{This ``1 versus 3 asymmetry'' is the origin of the structure in
four-dimensional spacetime consisting of one time axis and three spatial
axes.}
\end{quote}

In the Minkowski metric \(ds^2 = -dt^2 + dx^2 + dy^2 + dz^2\), only the
temporal component carries a negative sign. In the present model, this
sign asymmetry is \textbf{naturally derived} from the foundational fact
that ``exactly one axis was reversed to create negative volume.'' Rather
than inserting a sign into the metric after the fact, the sign asymmetry
originates from the very structure of the beginning of the universe.

This uniqueness does not hold in a three-dimensional cube model. In
three dimensions, \((-1)^1 = -1\) and \((-1)^3 = -1\), so both
single-axis reversal and triple-axis reversal produce a negative value.
There is no way to decide between ``1 versus 2'' and ``3 versus 0,'' and
the special status of the time axis cannot be uniquely determined. The
even-dimensional nature of four dimensions guarantees the structure in
which single-axis reversal uniquely determines the time axis.

\begin{center}\rule{0.5\linewidth}{0.5pt}\end{center}

\subsection{14. The Volume Conservation Law and the Beginning of the
Universe}\label{the-volume-conservation-law-and-the-beginning-of-the-universe}

The logical structure of the model is reorganized as follows:

\begin{enumerate}
\def\labelenumi{\arabic{enumi}.}
\tightlist
\item
  The volume conservation law cannot be violated (a fundamental
  principle of physics).
\item
  → If there was a beginning, the total volume must always be zero.
\item
  → Negative volume is logically required.
\item
  → Axis reversal is necessary.
\item
  → Pair symmetry cannot be maintained unless the dimension is even.
\item
  → Stability requires the integrality of \(\sqrt{n}\) so that corner
  cubes lie on integer-distance lattice points on the sphere.
\item
  → The smallest non-trivial solution is \(n = 4\).
\item
  → The minimal operation for negative volume is single-axis reversal.
\item
  → Exactly one axis is special = time.
\item
  → The remaining three axes are on equal footing = space.
\end{enumerate}

\textbf{There are only two assumptions:}

\begin{itemize}
\tightlist
\item
  \textbf{Assumption 1}: The universe is discrete.
\item
  \textbf{Assumption 2}: A stable minimal unit exists.
\end{itemize}

Four-dimensionality, the special nature of the time axis, and the
beginning of the universe are all \textbf{consequences} of these two
assumptions --- not input parameters.

\begin{center}\rule{0.5\linewidth}{0.5pt}\end{center}

\subsection{15. The Spherical Shell Model: Spacetime with Positive
Curvature}\label{the-spherical-shell-model-spacetime-with-positive-curvature}

Here one rule is changed. Instead of ``reducing the packing density
(spreading uniformly in the interior),'' consider the rule
``concentrating at the boundary of the dimension.'' What happens then?

The surface of a four-dimensional hypersphere is the three-dimensional
sphere \(S^3\). If the positive--negative pairs of hyperrectangles
concentrate on the surface rather than spreading uniformly in the
interior, the \(N(k)\) pairs adhere to the surface of a hypersphere of
radius \((2k+1)\) as an extremely thin shell.

All stability conditions are satisfied:

\begin{itemize}
\tightlist
\item
  The integrality \(\sqrt{4} = 2\) is maintained.
\item
  \(B_4\) symmetry is also maintained on the spherical shell.
\item
  The circumscribed sphere diameter \(2(2k+1)\) is unchanged.
\item
  The total volume zero of positive--negative pairs is preserved.
\end{itemize}

The decisive importance of this modification is that it naturally
explains \textbf{the origin of the three-dimensional space we observe}.
Since the surface of a four-dimensional hypersphere is a
three-dimensional manifold, an observer living on this thin shell would
perceive the universe as three-dimensional space.

This picture serves as a metaphor for \textbf{closed four-dimensional
spacetime with positive curvature}. In a closed Friedmann universe in
general relativity (FLRW metric with \(k = +1\)), the spatial
cross-section at each moment in time is \(S^3\), and its radius changes
with time. In the present model:

\begin{itemize}
\tightlist
\item
  Radius of \(S^3\) \(= (2k+1)\) → corresponds to time evolution
\item
  Thickness of shell \(\approx 0\) → matter concentrates on the spatial
  cross-section
\item
  Positive curvature \(= 1/(2k+1)^2\) → curvature decreases as \(k\)
  increases
\end{itemize}

As \(k\) increases, the curvature \(1/(2k+1)^2\) decreases, and the
space appears nearly flat locally. This is consistent with the
observational fact that the present universe appears nearly flat while
potentially retaining a small positive curvature.

\begin{center}\rule{0.5\linewidth}{0.5pt}\end{center}

\subsection{16. Universe and Anti-Universe: Separation of Positive and
Negative}\label{universe-and-anti-universe-separation-of-positive-and-negative}

A further rule is introduced:

\begin{quote}
\textbf{There may be multiple circumscribed hyperspheres, so long as
volume is conserved across the whole.}
\end{quote}

This relaxation of the constraint removes the requirement that positive
and negative hyperrectangles be packed into the same hypersphere. The
consequence is:

\begin{itemize}
\tightlist
\item
  \textbf{Positive hypersphere}: diameter \(2(2k+1)\), containing
  \(N(k)\) hyperrectangles with volume \(+1\)
\item
  \textbf{Negative hypersphere}: diameter \(2(2k+1)\), containing
  \(N(k)\) hyperrectangles with volume \(-1\)
\item
  \textbf{Total volume}: \(+N(k) - N(k) = 0\) (conservation law
  satisfied)
\end{itemize}

Both have the same diameter, the same count, and the same symmetry, and
each independently satisfies the stability conditions.

This signifies \textbf{pair creation of a universe and an
anti-universe}. The unfolding from a zero-dimensional state is realized
as twin universes with completely separated positive and negative
volumes.

\begin{center}\rule{0.5\linewidth}{0.5pt}\end{center}

\subsection{17. A Universe with Time and a Universe without
Time}\label{a-universe-with-time-and-a-universe-without-time}

Returning to the earlier argument that single-axis reversal generates
time:

\textbf{Our universe (the negative-volume side)}: - One axis reversed →
\((-1)^1 = -1\) → negative volume - A time axis exists → dynamic,
subject to change - Minkowski metric \((-,+,+,+)\)

\textbf{The paired universe (the positive-volume side)}: - No axis
reversal → \((+1)^4 = +1\) → positive volume - No time axis → all four
axes are on equal footing - Euclidean metric \((+,+,+,+)\)

Because the paired universe has no time, it has no change, causality, or
dynamics. It is a structure that simply ``exists'' statically as
four-dimensional space. Observing it is in principle impossible for us,
because the act of observation itself presupposes time.

\begin{quote}
\textbf{The existence of time is itself a ``liability,'' and a ``static
structure without time'' that cancels it must necessarily exist as a
pair.}
\end{quote}

Time is not something that exists as a matter of course; it is a
structure that arose as the price to be paid for satisfying the
conservation law.

\begin{center}\rule{0.5\linewidth}{0.5pt}\end{center}

\subsection{18. Fluctuations in Initial Distribution and the Complexity
of the
Universe}\label{fluctuations-in-initial-distribution-and-the-complexity-of-the-universe}

Consider initial conditions in which the distribution of positive and
negative hyperrectangles between the two hyperspheres is not perfectly
symmetric.

If the \textbf{positive hypersphere} contains \(+N\) and \(-m\) elements
while the \textbf{negative hypersphere} contains \(-N\) and \(+m\)
elements (with the total always summing to zero), each hypersphere is in
an unstable state in which a minority of opposite-sign elements is
present.

As this minority annihilates with the majority and rearrangement
proceeds, the following phenomena naturally emerge:

\textbf{Annihilation dynamics}: Positive-volume hyperrectangles that
have mixed into our hypersphere (the negative-volume side) pair up with
negative ones and disappear. This corresponds to the annihilation of
matter and antimatter (used here metaphorically) in the early universe.

\textbf{Residual asymmetry}: If the distribution is not perfectly
symmetric but has a minute bias, annihilation leaves a small
``remainder.'' This corresponds to baryon asymmetry (the reason only
matter remains in the present universe). Observations suggest that
matter outnumbered antimatter by approximately one part in a billion. In
the present model, this is expressed as a minute deviation in the
initial distribution.

\textbf{Seeds of structure formation}: Depending on the pattern of the
distributional deviation, positive and negative elements may fail to
fully separate locally and instead form clumps. These could become
density fluctuations and, in turn, seeds for galaxies and large-scale
structure.

\textbf{The complexity parameter}: When \(k\) is small, the total number
of hyperrectangles is small and the system reaches stability relatively
quickly. When \(k\) is large, the combinations are vast and a complex,
rich intermediate process emerges. The parameter \(k\) thus determines
the ``complexity'' of the universe.

\begin{center}\rule{0.5\linewidth}{0.5pt}\end{center}

\subsection{19. The Emergence of Time: Initially There Was No
Time}\label{the-emergence-of-time-initially-there-was-no-time}

Finally, the deepest consequence is stated.

Single-axis reversal is not the only way to realize negative volume.
Triple-axis reversal also gives \((-1)^3 = -1\), producing a negative
value. Consider an initial condition in which the reversal axes of the
negative-volume hyperrectangles are all different from one another.

\textbf{Initial state}: Some hyperrectangles have the first axis
reversed, others the third, and still others three axes simultaneously.
``Which axis is time'' is not determined. Every direction looks
equivalent, and \textbf{there is no consistent time axis --- that is,
there is no time}.

\textbf{Intermediate process}: As annihilation and rearrangement
proceed, the energetically most economical configuration (single-axis
reversal only) becomes dominant. This is \textbf{spontaneous symmetry
breaking}: from a state in which all four axes were equivalent, one is
progressively selected.

\textbf{Stable state}: Nearly all hyperrectangles come to share the same
single axis as their reversal axis. At this moment, for the first time a
consistent temporal direction is defined throughout the universe. This
is the \textbf{emergence of time}.

\begin{quote}
\textbf{Time is not a fundamental structure of the universe; it is an
order that emerged as a result of stabilization.}
\end{quote}

Initially there was no time and the four-dimensional space was
isotropic. Through annihilation dynamics the reversal axes aligned, and
time condensed.

This picture has contact with the following problems in physics:

\begin{itemize}
\tightlist
\item
  \textbf{The Planck era}: In the very early universe, the very concept
  of spacetime is thought not to be well-defined. In the present model,
  this is naturally expressed as a state in which the reversal axes are
  not aligned and time cannot be defined.
\item
  \textbf{Inflation}: The process by which the reversal axes rapidly
  align may correspond to exponential expansion. ``The birth of time''
  and ``the rapid expansion of space'' become two aspects of the same
  process.
\item
  \textbf{Contact with quantum gravity}: In many quantum gravity
  theories, ``time is not fundamental but emergent.'' The present model
  concretely realizes this as the statistical alignment of reversal axes
  of discrete hyperrectangles.
\end{itemize}

\begin{center}\rule{0.5\linewidth}{0.5pt}\end{center}

\subsection{20. Conclusion: What Follows from Two
Assumptions}\label{conclusion-what-follows-from-two-assumptions}

Counting the assumptions of this model precisely:

\begin{itemize}
\tightlist
\item
  \textbf{Assumption 1}: The universe is discrete.
\item
  \textbf{Assumption 2}: A stable minimal unit exists.
\end{itemize}

From these, via the volume conservation law, the following are all
\textbf{deductively derived}:

\begin{enumerate}
\def\labelenumi{\arabic{enumi}.}
\tightlist
\item
  \textbf{Spacetime is four-dimensional} --- the smallest non-trivial
  dimension for which \(\sqrt{n}\) is an integer
\item
  \textbf{Exactly one time axis is special} --- single-axis reversal as
  the minimal operation for negative volume
\item
  \textbf{Space is three-dimensional and isotropic} --- the remaining
  three axes, uninvolved in reversal, are on equal footing
\item
  \textbf{The sign structure of the Minkowski metric} \((-,+,+,+)\) ---
  the sign of the reversal axis
\item
  \textbf{The universe was able to ``begin''} --- expansion from volume
  zero without violating the conservation law
\item
  \textbf{The paired structure of universe and anti-universe} ---
  stabilization through separation of positive and negative
\item
  \textbf{Matter--antimatter asymmetry} --- a minute deviation in the
  initial distribution
\item
  \textbf{The emergence of time} --- statistical alignment of reversal
  axes
\item
  \textbf{The possibility of a closed universe with positive curvature}
  --- correspondence with the spherical shell model
\end{enumerate}

Even the number of dimensions is a derived quantity, not an input
parameter.

The present model requires no additional physical constants, no
compactification of extra dimensions, and no anthropic principle. From
two assumptions --- discreteness and stability --- and one principle ---
the volume conservation law --- it provides a unified framework for
addressing the principal questions concerning the structure of spacetime
and the origin of the universe.

\begin{center}\rule{0.5\linewidth}{0.5pt}\end{center}

\emph{This paper is a record of a model built up step by step through an
interactive thought experiment, starting from the elementary geometric
question of the circumscribed sphere diameter of a four-dimensional
hyperrectangle.}

\begin{center}\rule{0.5\linewidth}{0.5pt}\end{center}

\subsection{Appendix: Suggested Logical Reading
Order}\label{appendix-suggested-logical-reading-order}

The papers in this series are numbered by publication order. For logical
reading, the following order is recommended:

\begin{enumerate}
\def\labelenumi{\arabic{enumi}.}
\tightlist
\item
  Papers {[}1{]}--{[}3{]}: Foundational framework (central projection,
  symmetries, relativity of observation)
\item
  Paper {[}4{]}: \(R\) limits and nested structure
\item
  Paper {[}5{]}: Dimension addition
\item
  Paper {[}9{]}: Schwarzschild--de Sitter exact solution of Paper
  {[}1{]}'s equation and horizon structure at \(R \to 0\)
\item
  Paper {[}7{]}: Bridging geodesics
\item
  Paper {[}10{]}: Dimensional interpretation of the \(R \to 0\) limit
  (integration of nested structure, dimensional hierarchy, and exact
  solution)
\item
  Paper {[}6{]}: Subjective spaces at each dimension
\item
  \textbf{Paper {[}8{]}: Dimensional expansion from zero dimensions}
\end{enumerate}

\begin{center}\rule{0.5\linewidth}{0.5pt}\end{center}

\end{document}
